% =====================================================================
% MID SUBMISSION -- 2 October 2026.
% 7-8 pages excluding references, plus one optional appendix page.
% Must cover: a refined statement of the problem, the progress made, a
% concrete timeline for the rest of the project, and a literature review
% that establishes viability and motivation.
% Graded on the write-up, the code and a viva. Slides optional.
% Submit as ANLPDoomsday-Mid.zip (this PDF + the code).
% =====================================================================
\documentclass[11pt]{article}
\usepackage{acl}
\usepackage{times}
\usepackage{latexsym}
\usepackage[T1]{fontenc}
\usepackage[utf8]{inputenc}
\usepackage{microtype}
\usepackage{amsmath,amssymb}
\usepackage{booktabs}
\usepackage{enumitem}

\setlist{nosep,leftmargin=*}
\newcommand{\idf}{\mathrm{idf}}
\newcommand{\sfreq}{\mathrm{sf}}

\title{Predicting When Collection Statistics Help Top-$k$ Selection:\\
From Retrieval to Attention and Expert Routing}

\author{
  \textbf{Team ANLP Doomsday} \\[2pt]
  Aviral Gupta \quad Mohit Kumar Singh \quad Anirudh Sankar \quad Arihant Tripathy \\
  \texttt{\{aviral.gupta, mohit.singh, anirudh.sankar, arihant.tripathy\}@research.iiit.ac.in} \\
  IIIT Hyderabad, India
}

\begin{document}
\maketitle

\begin{abstract}
Retrieval, long-context attention and mixture-of-experts routing each select a few
items by an inner product against learned features and a hard top-$k$. Only retrieval
corrects that score, in BM25, with inverse document frequency, saturation and length
normalisation. We ask when the correction helps, and whether that can be predicted
before running it. In the completed retrieval setting, fourteen sparse representations
built from one frozen encoder are scored under all eight combinations of IDF, saturation and length normalisation,
in $182$ primary rows on thirteen collections and $1{,}226$ sensitivity rows on ten,
from two runs that agree to a median $0.0007$ nDCG@10. The correction helps in $90\%$
of primary rows. A Shapley decomposition assigns $61\%$ of the gain to saturation,
which needs no collection statistics, and $39\%$ to IDF and length together, within a
point on both runs. The
pre-registered predictor, a ridge regression on three frequency diagnostics, does no
better than the majority sign, because the effect of saturation depends on how unequal
the stored values are, which the index shows, while the value of IDF and length depends
on relevance, which it does not. Scoring each configuration
against a reference ranking supplies that information without labels. It predicts the
gain from collection statistics on unseen corpora with $R^2$ of $0.52$ to $0.66$, and
using collection statistics only where it predicts a gain recovers $61\%$ of the gap
to an oracle. All six pre-registered tests of the mechanism pass, five of them on data
the exploration had not used. Attention and routing are implemented and are next.
\end{abstract}

\section{Problem Statement}
\label{sec:problem}

\paragraph{Selection by inner product.} A learned sparse retriever scores a query $q$
against a document $d$ as $\langle u_q, u_d\rangle$ over a vocabulary of $V$
non-negative features and returns the top $k$. A block-sparse attention kernel scores
KV page $i$ for query $q_t$ as $\langle q_t, c_i\rangle$ over page summaries and
attends to the top-$k$ pages \citep{tang2024quest}. A mixture-of-experts router gates
on $\langle h_t, w_e\rangle$ and sends the token to its top-$k$ experts
\citep{shazeer2017moe,fedus2022switch}. All three compute an inner product and keep
the top $k$.

\paragraph{The correction.} For query and document vectors $x_q, x_d$ with
non-negative coordinates we score
\begin{equation}
S(q,d) = \sum_{j} x_{qj}\, w_j\, T(x_{dj}, \ell_d),
\label{eq:scorer}
\end{equation}
where $w_j$ is one or the smoothed inverse document frequency
$\idf_j = \ln(1 + (N - df_j + 0.5)/(df_j + 0.5))$, with $N$ documents of which $df_j$
have feature $j$, the length factor is
$\ell_d = 1$ or $1 - b + b L_d/\overline{L}$ with $L_d = \sum_j x_{dj}$, and
$T(f,\ell)$ is $f/\ell$ or the saturated $(k_1+1)f/(f + k_1\ell)$. Each of the three
components is a switch, so a configuration is a bit string $c \in \{0,1\}^3$ in the
order IDF, saturation, length. Configuration \texttt{000} is the plain inner product
and \texttt{111} is BM25 \citep{robertson2009bm25}. IDF and length normalisation read
collection statistics. Saturation reads only the document.

\paragraph{The unit of observation.} Collection frequency belongs to a corpus, so one
row of the analysis pairs a representation with a collection, at a support budget
(query and document feature caps), BM25 parameters $(k_1, b)$ and a basis seed. Each
row carries eight nDCG@10 values $Q_c$, one per configuration.

\paragraph{What is predicted.} The pre-registered target is the margin
$m = (Q_{111} - Q_{000})/\operatorname{sd}(Q_{000}, \ldots, Q_{111})$, a standardised
effect on the same scale in every setting. Section~\ref{sec:progress} shows that
$m$ is poorly shaped for prediction, and we add targets defined on the same eight
values. The Shapley value $\phi_s$ of switch $s$ is its marginal effect
averaged over every order in which the switches can be turned on
\citep{shapley1953value}. Writing $Q_A$ for the configuration with the switches in $A$
on,
\begin{equation}
\phi_s = \sum_{A \not\ni s}
\tfrac{|A|!\,(2-|A|)!}{6}\bigl(Q_{A\cup\{s\}} - Q_{A}\bigr),
\label{eq:shapley}
\end{equation}
and the three values sum exactly to $Q_{111} - Q_{000}$. The collection gain
$G = Q_{111} - Q_{010}$ is the value of IDF and length normalisation once saturation
is applied. Its sign is the decision the project is about, whether to use collection
statistics for a given representation and collection. A predictor is a function of
label-free quantities that returns $m$, $\phi_s$ or $G$ for a row. It is scored by
sign accuracy against the majority sign, held-out $R^2$ and Spearman correlation, and,
for $G$, by the nDCG@10 obtained when collection statistics are used only where the
prediction is positive.

\paragraph{Transfer.} Eq.~\eqref{eq:scorer} applies to any selection with a
non-negative score and a selection count. For KV page $i$ at decode step $t$, the
feature value is the softplus of the Quest criticality $r_{t,i}$, the collection
frequency is $\sfreq_i$, the fraction of calibration queries for which page $i$ enters
the top-$k$ among the pages that exist at that point, and the length is the norm of the
page's concatenated key minima and maxima. For expert $e$ at a token, the value is the
router probability $p_e$, the frequency is the share of calibration tokens routed to
$e$, and the length is the norm of the router logits. The selected experts are combined
with their original gate weights, so the correction changes only which experts run.
In attention a row is a model, RULER task, context length and token budget, with RULER accuracy as the metric; in routing it is a model and task, with
negative token log-likelihood on WikiText and zero-shot MMLU accuracy. A predictor
fitted in retrieval is applied there without refitting. It transfers if its predicted
sign is correct on at least $70\%$ of rows and its mean absolute error is at most $1.5$
times its corpus-held-out RMSE in retrieval.

\section{Motivation and Research Gap}
\label{sec:motivation}

The inner-product selectors of Section~\ref{sec:problem} use the raw match score,
with no account of how often a feature is selected across the collection, how
evidence accumulates within an item, or how large the item is. Classical lexical
retrieval corrects its raw match score for all three, with parts that arrived separately, inverse document frequency \citep{sparckjones1972idf}, term-frequency
saturation \citep{robertson1994poisson} and document length normalisation
\citep{singhal1996pivoted}, combined in BM25 \citep{robertson2009bm25}. Inside the
transformer only fragments have been reinvented. Loss-free load balancing adds a
per-expert bias driven by recent load \citep{wang2024lossfree}, an additive frequency
discount with no saturation or length term, and streaming inference keeps the
most-attended early keys as an exception \citep{xiao2024streamingllm}.

Three papers have scored the same learned features both with and without the
correction, all in retrieval. BM25 wins on sparse-autoencoder features of a dense
retriever \citep{clavie2026latent} and on SAE concept units \citep{park2025clsr}, and
comes out level on visual words \citep{han2026bm25v}. Clavi\'e et al.\ and Han et al.\
give the same explanation, that BM25 helps while a vocabulary keeps a Zipfian frequency
curve. None of the three varies the vocabulary systematically, and none attributes the gain to the
three components. The Zipfian account has not been turned into a quantity that can be
computed before the experiment, nor checked where the same computation runs inside a
model.

\section{Contributions}
\label{sec:contributions}

\begin{enumerate}
\item A controlled grid, built and run. Fourteen sparse representations from one
frozen encoder, seven bases at two granularities, are scored under all eight
configurations on thirteen collections, with sensitivity families over support
budget, $k_1$ and $b$, basis seed and collection-statistic controls, in 182 primary
rows and 1{,}226 sensitivity rows from two independent runs.
\item An attribution of the gain. Shapley values over the three switches show that
saturation, a per-document transform, carries $61\%$ of the gain and collection
statistics $39\%$, and that IDF alone carries about a tenth, falling as the support
budget grows.
\item A negative result for the pre-registered predictor, with its cause. Frequency
diagnostics predict neither the margin nor the collection gain, because whether IDF
and length help depends on relevance, which no statistic of the index measures.
\item A label-free predictor that works in retrieval. Agreement of each configuration
with a reference ranking predicts all three components and the collection gain on
held-out basis families and corpora. Five pre-registered tests of the mechanism pass on
data the exploration did not use.
\item A transfer test outside retrieval, implemented. KV-page selection and expert
routing adapters with native-parity checks are ready for the predictor to be frozen
and applied.
\end{enumerate}

\section{Background}
\label{sec:background}

The correction needs three things of a representation, non-negative coordinates, a
meaningful per-feature collection frequency $df_j$, and a magnitude that admits a
length. Dense embeddings lack the first two, since every coordinate is signed and
active in every document, so every $df_j$ equals $N$. Sparse codes built from those embeddings have all three, and
so does a selection count over KV pages or experts.

For a fixed linear map $C$ and diagonal weights $D$, weighted feature scoring is
bilinear in the original vectors, $S(q,d) = (Cq)^\top D (Cd) = q^\top C^\top D C d$,
and it preserves every ranking of the original inner product when
$C^\top D C = \alpha I$ with $\alpha > 0$. Split rectification, support truncation,
saturation and a document-dependent length factor are not linear, so none of the
representations here meets the condition, and the grid is needed to learn what the
correction does. The identity, the sufficient conditions and a counterexample in which
a diagonal rescaling reverses a ranking are checked numerically in the code.

\section{Related Work}
\label{sec:related}

\paragraph{Term weighting.} IDF was proposed as a statistical account of term
specificity \citep{sparckjones1972idf}. Saturation came from approximating the
two-Poisson model of within-document frequency \citep{robertson1994poisson}, and
length normalisation was pivoted so that long documents are neither favoured nor
penalised \citep{singhal1996pivoted}. BM25 combines the three
\citep{robertson2009bm25}. Each part was motivated by word statistics, which is why
it is open whether it helps features that are not words.

\paragraph{Retrieval scorers.} Dense retrievers score by the inner
product of embeddings \citep{karpukhin2020dpr}; the encoder used here is multilingual
E5 \citep{wang2024me5}. Learned sparse retrieval keeps the inverted index and learns
term weights, through contextual exact match \citep{gao2021coil} or expansion over the
subword vocabulary in SPLADE \citep{formal2021splade,formal2021spladev2,lassance2024spladev3}.
\citet{lin2021conceptual} and \citet{nguyen2023unified} place these methods in one
framework, and \citet{mackenzie2021wacky} show that learned weights change the cost of
query processing, which is why we report posting visits and index size. All of them
score by inner product.

\paragraph{Learned concept vocabularies.} Sparse autoencoders decompose dense
activations into many non-negative features \citep{cunningham2024sae}, and Top-$K$
SAEs fix the number active per input \citep{gao2025topksae}. In retrieval, CL-SR indexes SAE-decomposed embeddings as concept units and weights them by IDF,
and dropping that weighting costs 11 to 18 percent of its nDCG@10 on TREC-DL
\citep{park2025clsr}, the one published ablation of the correction. SPLARE trains an
SAE-based sparse retriever \citep{formal2026splare} and SAE-SPLADE replaces SPLADE's
language-model head with an SAE \citep{zong2026tokens}; both score by inner product.
Quantised codes are the other route to a discrete vocabulary, through product
quantisation \citep{jegou2011pq}, residual quantisation \citep{chen2010rq} and their
GPU implementations \citep{johnson2021faiss}.

\paragraph{BM25 on learned features.} \citet{clavie2026latent} show that SAE features of
a frozen dense retriever have roughly Zipfian collection statistics and can be scored
by BM25 directly. In their Appendix~C, BM25 reaches $0.582$ nDCG@10 against $0.381$ for
the inner product on SAE latent terms and $0.455$ against $0.568$ on SPLADE-v3, so the
scorers swap places when the vocabulary changes. \citet{han2026bm25v} find the two
within about two points of R@1 on image retrieval and attribute the parity to top-$k$
filtering already suppressing what IDF would down-weight. Raising their per-patch
budget from $16$ to $128$ makes document frequencies converge and R@1 collapse from
$0.508$ to $0.038$. Our budget sweep shows the matching trend, the share of the gain
due to IDF falling as the budget grows (Section~\ref{sec:attribution}).

\paragraph{Predicting effectiveness without labels.} Query performance prediction
estimates retrieval quality from the query and the ranking alone
\citep{carmel2010difficulty}, through the divergence of the retrieved language model
from the collection \citep{cronen2002clarity} or the spread of the top retrieval
scores \citep{shtok2012nqc}. Retrieval systems can also be ranked without judgements
by comparing each system with pseudo-relevant documents drawn from the pooled output
\citep{soboroff2001ranking}. Our label-free predictor uses both ideas, score-shape
statistics and agreement with a reference ranking, to predict a difference between
scorers, which neither line of work targets.

\paragraph{Attributing an effect to components.} The Shapley value assigns a joint
outcome to players by averaging marginal contributions over orderings
\citep{shapley1953value}, and SHAP applies it to model predictions
\citep{lundberg2017shap}. We apply it to a complete $2^3$ factorial, where it is exact
and needs no sampling.

\paragraph{Selection inside the transformer.} Quest scores KV pages by an inner
product against page min-max key summaries \citep{tang2024quest}. H2O keeps recent and
heavy-hitter tokens \citep{zhang2023h2o}, SnapKV keeps the positions an observation
window attends to \citep{li2024snapkv}, StreamingLLM keeps the first keys as attention
sinks \citep{xiao2024streamingllm}, and Expected Attention prunes pairs that future
queries are predicted to ignore \citep{devoto2025expectedattention}. Massive activations
single out a few positions by magnitude \citep{sun2024massive}. None of these uses the
rate at which an entry is selected across queries. Long-context quality is measured by
RULER \citep{hsieh2024ruler}.

\paragraph{Expert routing.} Sparse mixture-of-experts layers route each token to its
top-$k$ experts \citep{shazeer2017moe,lepikhin2021gshard,fedus2022switch}. Balance is
enforced by an auxiliary loss \citep{zoph2022stmoe} or, without one, by a load-driven
bias \citep{wang2024lossfree}. DeepSeekMoE splits experts more finely and adds shared
experts to sharpen their specialisation \citep{dai2024deepseekmoe}. \citet{su2026superexperts} find a handful of experts
by activation magnitude whose removal degrades the model badly. Our routing setting
uses OLMoE \citep{muennighoff2025olmoe}, which activates 8 of 64 experts per layer. To our knowledge no work treats retrieval scoring, KV
selection and routing as one scoring problem.

\paragraph{Viability.} Two papers report the
correction winning on SAE features and one reports parity, all without a systematic
vocabulary variation or a component attribution. The correction is cheap to switch on,
sparse codes are cheap to build from one encoder, and label-free proxies for quality
already work for single systems, so closing the gap needs no new model.

\section{Method}
\label{sec:method}

\paragraph{Representations.} All representations come from
\texttt{multilingual-e5-base} \citep{wang2024me5}, frozen at a pinned revision. Seven
bases are fitted on the same $200{,}000$ MS~MARCO passages \citep{bajaj2016msmarco}
and nowhere else, namely split-rectified identity, Gaussian projection (768) and PCA
(384), soft $k$-means over $4{,}096$ centroids, product and residual quantisation
($8 \times 256$), and a Top-$K$ SAE with $8{,}192$ features of which at most 32 are
active. Each is applied at pooled
granularity, to the mean of the content token states, and at token granularity, to
each state before summing, giving fourteen representations. A secondary SAE shares one
dictionary across both granularities to separate aggregation from fitting. SPLADE
has its own backbone, so it is an external comparison outside the grid.

\paragraph{Corpus statistics.} Supports are capped before any statistic, at 32
query and 128 document features in the primary budget. Document frequency, lengths and
IDF are computed on each target corpus from the capped vectors, and retrieval is exact
over the whole corpus. The primary grid uses $k_1 = 1.2$, $b = 0.75$ and seed zero.

\paragraph{Collections.} The thirteen collections, from eleven corpora, are MS~MARCO
dev with TREC-DL 2019 and 2020 over the same documents, six BEIR sets
\citep{thakur2021beir}, LIMIT \citep{weller2025limit}, and the Hindi, Bengali and
Telugu portions of MIRACL \citep{zhang2023miracl}, each evaluated within its language.
LIMIT is built so that the top-$k$ sets its queries need exceed what a low-dimensional
embedding can represent.

\paragraph{Experiment families.} The primary family (Table~\ref{tab:families}) was run
twice, on a four-GPU cluster node for all thirteen
collections and on a single-GPU workstation for the ten outside MS~MARCO, which also
ran every sensitivity family.

\begin{table}[t]
\centering
\small
\begin{tabular}{@{}lr@{}}
\toprule
Family & Rows \\
\midrule
primary (14 representations $\times$ 13 collections) & 182 \\
support budget 16/64, 32/128, 64/256, native & 462 \\
$k_1 \in \{0.6, 1.2, 2.4\}$, $b \in \{0, 0.5, 0.75\}$ & 630 \\
basis seeds 1 and 2 & 100 \\
shuffled and uniform $df$ controls & 84 \\
shared SAE dictionary & 20 \\
\midrule
union on the ten non-MS~MARCO collections & 1{,}226 \\
\bottomrule
\end{tabular}
\caption{Retrieval rows, each with eight scoring configurations. The primary count is
the cluster run's; the other families and the union are the workstation run's. Families
share their default cells, so the union is smaller than the sum. The MS~MARCO budget rows are
pending (Section~\ref{sec:plan}).}
\label{tab:families}
\end{table}

\paragraph{Significance.} Comparisons use a paired bootstrap over queries that recomputes
numerator and denominator, with Holm correction \citep{holm1979} within each
pre-declared family, the all-on against all-off comparison and the three component main
effects. The cluster run drew $2{,}000$ samples, which puts the smallest attainable Holm
$p$-value over 182 comparisons at $0.091$, so significance is reported from the
workstation run, which drew $10{,}000$.

\paragraph{Diagnostics.} None of the features describing a row reads relevance labels.
The six pre-registered frequency diagnostics are the active fraction, the Zipf slope of the
feature-frequency curve, the Gini coefficient of activation mass, frequency and
activation mass in the most frequent $1\%$ of features, and IDF variance; the primary
ridge model uses the slope, IDF variance and frequent activation share. Index
statistics describe the stored values, chiefly $\mathrm{cv}_x$, the coefficient of
variation of the stored document values, with the peak-to-mean ratio and the spread of
document lengths. Ranking features are read from the stored top-100 list $R_c(q)$ of every
configuration $c$ and query $q$. The disruption of switch $s$ is
$1 - |R_c^{10}(q) \cap R_{c'}^{10}(q)|/10$ averaged over queries and over the four
pairs $(c, c')$ that differ only in $s$. Score-shape statistics of the kind used in
query performance prediction \citep{shtok2012nqc} are the coefficient of variation of
the top-100 scores and the relative drop from the first to the tenth. The
pseudo-quality $\tilde{Q}_c$ is the mean nDCG@10 of $R_c(q)$ when the top 10 of a
reference retriever on the same query is taken as relevant, with binary gains. The
references are text BM25, the dense encoder, and the two fused by reciprocal rank
($k = 60$ over their top 100). Replacing $Q_c$ by $\tilde{Q}_c$ in
Eq.~\eqref{eq:shapley} and in $G$ gives a pseudo-label estimate of every target,
$\tilde{\phi}_s$ and $\tilde{G}$, with no fitted parameter.

\paragraph{Validation.} Predictors are ridge regressions with the scaler and the
regularisation chosen inside training folds, and gradient-boosted trees with fixed
settings. Folds hold out a whole basis family (linear, clustered, quantised, learned) or
a whole corpus, reported separately, since a new basis and a new corpus are different
claims. A fit on primary rows is also scored on the budget, seed and parameter rows,
which share collections with it.

\paragraph{Transfer adapters.} The attention adapter implements Quest-style page
criticality over sixteen-token pages with dense prefill and sparse decode, counts how
often each page enters the top-$k$ over causal prefill positions, and applies
Eq.~\eqref{eq:scorer} with pages for documents; all-off calls the native selector, and
an H2O reference shares the budget accounting. The routing adapter applies the same
correction to OLMoE router probabilities for selection while combining with the
original gate weights; all-off calls the native router. Six tests confirm exact
native parity, causal calibration, budget charging including sink pages, and the gate
mixture.

\section{Progress}
\label{sec:progress}

\subsection{The pipeline reproduces itself}

The two runs used the same imported bases and pinned data. Over the 140 primary rows
they share, the median absolute difference in all-on nDCG@10 is $0.0007$ (maximum
$0.0077$, from fp16 rounding on different GPUs), and the margin sign agrees in 136 of
the 137 rows where it is defined. The text BM25 baselines agree to seven decimals. Basis
seeds move all-on nDCG@10 by a median standard deviation of $0.006$ over 50
representation-collection groups. The conceptual replication of the Appendix~C reversal
of \citet{clavie2026latent} needs SPLADE-v3 and is scheduled (Section~\ref{sec:plan}).

\subsection{Retrieval quality}

\begin{table}[t]
\centering
\small
\setlength{\tabcolsep}{4pt}
\begin{tabular}{@{}lrrrrr@{}}
\toprule
Collection & BM25 & E5 & \texttt{000} & \texttt{111} & best \texttt{111} \\
\midrule
MS MARCO dev & .222 & .425 & .013 & .130 & .318 \\
TREC-DL 19 & .456 & .704 & .036 & .255 & .604 \\
TREC-DL 20 & .484 & .697 & .035 & .251 & .556 \\
ArguAna & .411 & .445 & .021 & .130 & .371 \\
FiQA & .236 & .382 & .013 & .085 & .266 \\
NFCorpus & .307 & .325 & .064 & .115 & .285 \\
SCIDOCS & .152 & .172 & .009 & .044 & .120 \\
SciFact & .662 & .694 & .089 & .197 & .419 \\
TREC-COVID & .582 & .695 & .134 & .246 & .508 \\
LIMIT & .939 & .016 & .001 & .004 & .032 \\
MIRACL hi & .357 & .583 & .011 & .124 & .328 \\
MIRACL bn & .407 & .701 & .025 & .216 & .490 \\
MIRACL te & .331 & .751 & .019 & .259 & .579 \\
\bottomrule
\end{tabular}
\caption{nDCG@10 of text BM25, the dense encoder, and the fourteen sparse
representations with the plain inner product (\texttt{000}) and full correction
(\texttt{111}), averaged, with the best representation's \texttt{111}. The best is the
token-level SAE everywhere except LIMIT (token $k$-means).}
\label{tab:quality}
\end{table}

In the primary grid (Table~\ref{tab:quality}), the plain inner product of the capped
sparse codes nearly fails, with mean nDCG@10 below $0.14$ on every collection, and full
correction raises the mean by a factor of between $1.8$ and $14$. Token granularity is
far stronger than pooled (mean all-on $0.219$ against $0.097$), and the token-level SAE is the best
representation on twelve collections, reaching $0.60$ on TREC-DL 2019 against $0.70$
for the dense encoder it was built from. On the ten workstation collections the pooled
SAE reaches $0.095$ with its own dictionary and $0.291$ with the token-trained one,
against $0.339$ at token level, so most of the pooled SAE's deficit comes from fitting
its dictionary on pooled vectors, and aggregation order accounts for the rest. LIMIT behaves as its construction predicts, BM25 on words reaching
$0.94$ where the dense encoder and every representation built from it stay near zero.

\subsection{Whether the correction helps}

The margin is positive in $90\%$ of the 179 primary rows with a defined margin. Over
the 140 rows of the workstation run, Holm-corrected all-on against all-off comparisons
are significant in 99, and the main effects of saturation, length and IDF in 93, 79 and
56, with mean effects of $+0.065$, $+0.032$ and $+0.013$ nDCG@10. Of the 182 primary
rows, full correction is the best of the eight configurations in only 43; the most
frequent winner, in 51, is saturation with length normalisation and no IDF. On the
workstation run the all-off configuration is below $0.01$ nDCG@10 in $42\%$ of rows,
and the margin falls between $1.5$ and $3.5$ in $83\%$, so the margin mostly records
that the uncorrected scorer is broken and leaves little variation to predict.

\subsection{Where the gain comes from}
\label{sec:attribution}

\begin{table}[t]
\centering
\small
\setlength{\tabcolsep}{4pt}
\begin{tabular}{@{}lrrrr@{}}
\toprule
Rows & IDF & Sat. & Len. & IDF + len. \\
\midrule
cluster, 182 primary & .10 & .62 & .28 & .38 (.36--.41) \\
workstation, 140 primary & .10 & .61 & .29 & .39 (.36--.42) \\
\midrule
budget 16/64 & .16 & .60 & .25 & .40 \\
budget 64/256 & .08 & .57 & .35 & .43 \\
native, 3 collections & $-.02$ & .55 & .47 & .45 \\
\bottomrule
\end{tabular}
\caption{Share of the all-on minus all-off gain assigned to each switch by
Eq.~\eqref{eq:shapley}. Intervals resample whole collections. The primary budget is
32/128; the other budget rows cover the same ten collections as the workstation primary
row, and native support covers ArguAna, NFCorpus and SciFact.}
\label{tab:shapley}
\end{table}

Saturation carries about $61\%$ of the gain and IDF about $10\%$ on both runs
(Table~\ref{tab:shapley}). Once saturation is applied, collection
statistics add $0.015$ nDCG@10 on average on the workstation run and $0.021$ on the
cluster run, and help in about two thirds of rows. On the ten workstation collections,
the share of IDF falls from $16\%$ at the 16/64 budget to $8\%$ at 64/256 while that of
length rises from $25\%$ to $35\%$, and at native support on three collections IDF's
share is below zero. \citet{han2026bm25v} report the same trend for image features and
attribute it to document frequencies converging as budgets grow. On the five
collections of the parameter family, all-on nDCG@10 at $b = 0$ falls as $k_1$ grows
($0.170$, $0.164$, $0.153$), and at $b = 0.75$ it is flat ($0.175$ to $0.176$), so weaker
saturation costs most when length normalisation is off.

Replacing document frequencies by shuffled ones on
three collections turns the mean IDF Shapley value from $+0.011$ to $-0.011$ (real
larger in $83\%$ of 42 rows, one-sided Wilcoxon $p = 8.9\times10^{-7}$); uniform
frequencies remove it by construction, and do not lower all-on nDCG@10 ($0.157$
against $0.147$, $p = 0.15$). Real frequencies carry information, but little of it shows in
nDCG@10.

\begin{table}[t]
\centering
\small
\setlength{\tabcolsep}{3.5pt}
\begin{tabular}{@{}llrrrrr@{}}
\toprule
Basis & Gran. & \texttt{111} & $\mathrm{cv}_x$ & $\phi_{\mathrm{sat}}$ & IDF + len. & $G > 0$ \\
\midrule
identity & pooled & .147 & 0.37 & .043 & .68 & 1.00 \\
 & token & .142 & 0.60 & .060 & .55 & .77 \\
Gaussian & pooled & .098 & 0.29 & .020 & .77 & .85 \\
 & token & .088 & 0.55 & .029 & .64 & .92 \\
PCA & pooled & .222 & 0.75 & .138 & .24 & .69 \\
 & token & .311 & 1.35 & .155 & .45 & .46 \\
$k$-means & pooled & .018 & 0.04 & .000 & .94 & .31 \\
 & token & .185 & 1.71 & .125 & .11 & .85 \\
PQ & pooled & .012 & 0.00 & .000 & 1.00 & .62 \\
 & token & .144 & 1.44 & .099 & .09 & .69 \\
RQ & pooled & .056 & 0.00 & .000 & 1.00 & .92 \\
 & token & .292 & 2.10 & .176 & .18 & .69 \\
SAE & pooled & .127 & 0.18 & .014 & .66 & .77 \\
 & token & .373 & 3.31 & .193 & .32 & .46 \\
\bottomrule
\end{tabular}
\caption{Per representation, mean all-on nDCG@10, mean saturation Shapley value, share
of the gain from IDF and length, and the fraction of collections where collection
statistics help after saturation ($G > 0$), over the 13 primary collections, with the
mean $\mathrm{cv}_x$ of the stored values over the ten whose indexes were measured.}
\label{tab:bases}
\end{table}

Across representations (Table~\ref{tab:bases}), saturation's contribution rises with
$\mathrm{cv}_x$, and the extreme case is exact. Pooled PQ
and RQ codes store equal values with equal lengths ($\mathrm{cv}_x = 0$), and their
Shapley values for saturation and length are exactly zero in all 26 of their cluster
rows and all 166 workstation rows, as the mechanism below requires. The share of the
gain from collection statistics is largest for the weakest representations. For the two
strongest, token-level SAE and PCA, collection statistics help after saturation on
fewer than half the collections, which is where predicting the decision matters.

\subsection{The pre-registered predictor}

The primary ridge model predicts the margin with held-out $R^2$ of $0.04$ across basis
families and $0.02$ across corpora on the cluster run, and $-0.11$ and $0.14$ on the
workstation run with the rows that respond to fewer than two switches excluded. Its
sign accuracy equals the majority-sign baseline in every case ($0.899$ and $0.924$). The
six-diagnostic model reaches $R^2 = 0.30$ across corpora on the workstation run but
$-0.54$ across basis families, and its sign accuracy is never more than $0.03$ above the
majority sign. By the pre-registered criterion the frequency diagnostics carry no usable
information about the margin.

\subsection{What predicts the gain}

\begin{table}[t]
\centering
\small
\setlength{\tabcolsep}{3pt}
\begin{tabular}{@{}llrrr@{}}
\toprule
Target & Features & Family & Corpus & Params \\
\midrule
$\phi_{\mathrm{sat}}$ & pre-registered & $-.46$ & .14 & .02 \\
 & index & .47 & .42 & .50 \\
 & pseudo-labels & .56 & .62 & .76 \\
$\phi_{\mathrm{len}}$ & pre-registered & $-.30$ & .14 & .00 \\
 & index & $-.16$ & .25 & $-.02$ \\
 & pseudo-labels & .49 & .70 & .77 \\
$\phi_{\mathrm{idf}}$ & pre-registered & $-.01$ & .00 & .02 \\
 & index & .02 & .02 & .07 \\
 & pseudo-labels & .40 & .21 & .73 \\
$G$ & pre-registered & $-.07$ & $-.02$ & .16 \\
 & index & $-.05$ & $-.09$ & .15 \\
 & pseudo-labels & .38 & .52 & .62 \\
 & pseudo-labels, E5 & .51 & .66 & .80 \\
\bottomrule
\end{tabular}
\caption{Held-out $R^2$ of ridge regression fitted on 139 primary rows, holding out
basis families, corpora, or predicting the 560 rows at other $k_1$ and $b$. Each tier
adds features to the one above it. Pseudo-labels use text BM25 unless marked E5.}
\label{tab:predictor}
\end{table}

\paragraph{Mechanism.} With length normalisation off, saturation maps a stored value
$f$ to $g(f) = (k_1+1)f/(f+k_1)$, and $g(f)/f$ falls as $f$ grows, so it shrinks large
values relative to small ones. If every stored value were equal, $g$ would rescale all
scores alike and change no ranking; the change grows with how unequal the values are,
which $\mathrm{cv}_x$ measures, and a smaller $k_1$ bends $g$ more. How much saturation
changes the ranking is a property of the index, and the change helps in $91\%$ of
cluster primary rows outside the two equal-valued codes. Length normalisation and IDF change the
ranking by an amount the index fixes, but whether the change helps depends on whether
heavy documents are over-rewarded relative to their relevance, and whether rare features
separate relevant documents from the rest. Both are properties of the index jointly
with relevance, so no statistic of the index can fix their sign.

\paragraph{Prediction.} Index statistics predict saturation and fail for IDF and the
collection gain, and the frequency diagnostics predict nothing (Table~\ref{tab:predictor}).
Pseudo-labels carry the relevance information the index lacks and predict every
target. The Shapley
values and $G$ are fixed linear combinations of the eight $Q_c$, so a pseudo-quality
that orders configurations like the true one gives correlated targets, and it does, with
a median within-row Spearman correlation of $0.98$ between pseudo and true nDCG@10 over
the eight configurations of 1{,}112 rows. Used with no fitting, the pseudo-label
estimates rank the true values with Spearman $0.91$ for saturation, $0.82$ for length,
$0.83$ for IDF and $0.66$ for $G$ on the primary rows. Gradient-boosted trees do no
better on the family and corpus axes.

\paragraph{Decision.} The decision rule uses collection statistics only where the
unfitted pseudo-label estimate $\tilde{G}$ is positive, and has no fitted parameter. It
raises mean nDCG@10 on the primary rows from $0.1429$
(always) to $0.1478$, against $0.1510$ for an oracle, recovering $61\%$ of the headroom
with a collection-resampled $95\%$ interval of $0.0003$ to $0.0101$ above always. The
rule recovers $67\%$ at other budgets and $65\%$ at other parameters, and $22\%$ at other
seeds, where the interval reaches zero. Across all rows, the rule gains most on SciFact,
where always using collection statistics loses $0.012$ on average, and on NFCorpus,
where always using them gains nothing; it is at chance only on LIMIT and TREC-COVID, where the true
gain is near zero.

\paragraph{Choice of reference.} Agreement with the dense reference partly measures how
much of the embedding a configuration keeps, since the sparse codes are built from the
same encoder. Agreement with BM25 may be favoured by the judgements, since lexical
systems contributed to the pools behind several collections. The result holds with
either (the E5 row of Table~\ref{tab:predictor}), within-row agreement is $0.98$ for
both, and the unfitted rule with the dense reference recovers $45$ to $75\%$ of the
headroom across the four groups of rows. All of this analysis was chosen after the
pre-registered result and is exploratory.

\begin{table*}[t]
\centering
\small
\begin{tabular}{@{}llp{0.58\textwidth}l@{}}
\toprule
Week & Dates & Work & Owner \\
\midrule
1 & 3--9 Oct & Resume the MS~MARCO budget sweep (84 rows) after the cluster maintenance.
Run SPLADE-v3 as the external comparison and a conceptual replication of the Appendix~C
reversal. Commit the second pre-registration, the pseudo-label rule tested on the unseen
MS~MARCO budget rows, and the frozen transfer predictors. & Aviral, Mohit \\
 & & Validation runs of both adapters on the cluster, native parity at scale, calibration
cost and memory. & Anirudh, Arihant \\
2 & 10--16 Oct & Setting B primary grid on Llama-3.2-1B-Instruct, all 13 RULER tasks with
500 examples, 4K and 8K contexts, budgets of 512 and 1{,}024 tokens, with H2O. Second
confirmatory round in retrieval. & Anirudh, Mohit \\
3 & 17--23 Oct & Setting C on OLMoE-1B-7B, WikiText-2 negative log-likelihood and zero-shot
MMLU, with the additive loss-free-balancing analogue. Written proof and checks of the
bilinear identity and its ranking conditions. & Arihant, Mohit \\
4 & 24--28 Oct & Transfer scoring with uncertainty for both predictors, per setting.
Figures and tables regenerated from saved artifacts. & all \\
5 & 29--31 Oct & Final report, slides, code archive. & all \\
\bottomrule
\end{tabular}
\caption{Schedule to the final submission. If time runs short we cut, in order, GSM8K and
expert preservation, the 8K context in Setting B, MMLU in Setting C (keeping
WikiText), and the Appendix~C replication.}
\label{tab:plan}
\end{table*}

\subsection{Confirmatory tests}

Six hypotheses were committed to the repository before their tests ran, on the $k_1$ and
$b$ family, whose 630 rows the exploration up to that point had not used. H2 to H6 are
Holm-corrected together; H1 overlaps the exploratory data and checks consistency only.

\begin{table}[t]
\centering
\small
\setlength{\tabcolsep}{3pt}
\begin{tabular}{@{}llr@{}}
\toprule
 & Hypothesis & Result \\
\midrule
H1 & IDF + length share at default in $[.36, .42]$ & $.393$ \\
H2 & $\phi_{\mathrm{sat}}$ at $k_1 = 0.6$ exceeds $k_1 = 2.4$ & $.093 > .080$ \\
H3 & within-corpus $\rho(\mathrm{cv}_x, \phi_{\mathrm{sat}}) \ge 0.3$ & $.89$, $.89$ \\
H4 & $\mathrm{cv}_x$ model ranks $\phi_{\mathrm{sat}}$ at $\rho \ge 0.5$ & $.87$, $.90$ \\
H5 & $\phi_{\mathrm{len}} = 0$ exactly at $b = 0$ & max $0.0$ \\
H6 & $\phi_{\mathrm{len}}$ at $b = 0.75$ exceeds $b = 0.5$ & $.045 > .029$ \\
\bottomrule
\end{tabular}
\caption{Confirmatory tests, all passing. H2 and H6 have Holm-corrected
$p = 3.6\times10^{-5}$; H3 and H4 are tested at $k_1 = 0.6$ and $2.4$, and H4's model is fitted
on primary rows at $k_1 = 1.2$; H5 covers 210 rows.}
\label{tab:confirm}
\end{table}

All six pass (Table~\ref{tab:confirm}). H5 is a check of the implementation. H2, H3,
H4 and H6 confirm that the switches behave as their mechanisms predict when their strength is
varied, and that the value-dispersion predictor of saturation keeps its ranking at
settings it was not fitted on.

\subsection{Transfer settings}

The adapters and their parity tests are complete (Section~\ref{sec:method}); no
production run has been made. Both predictors and the acceptance thresholds will be
frozen before any corrected transfer outcome is seen.

\section{Remaining Plan}
\label{sec:plan}

The pre-registered predictor is kept and will be carried to attention and routing as registered. The pseudo-label predictor is
carried beside it, with a reference native to each setting, exact attention mass for KV
pages, and the dense softmax mixture over all experts for routing, both label-free and
computed from the uncorrected model. Before any transfer outcome is seen, both
predictors, the references and the thresholds are committed, in the first week of the
schedule in Table~\ref{tab:plan}.

\section{Limitations}
\label{sec:limitations}

All representations come from one encoder, so basis and encoder effects are not
separable, and E5 was trained with MIRACL data, so the Indic results do not concern
unseen languages. The sparse codes are far below the dense encoder they come from,
which bounds how far the attribution speaks to strong learned sparse retrievers; the
SPLADE comparison is meant to test that. The label-free predictor needs queries and a
second retriever run, and its evidence is exploratory until the second pre-registered
round. Each reference carries a bias, and the result holding with both makes a single
bias an unlikely explanation. Routing is inference-time only, so nothing is claimed about training
dynamics.

\bibliography{refs}

\end{document}
